\section{Reconstructed R1 unbounded-$q$ source and scalar certificate}
\label{sec:reconstructed-r1-qo-source-scalar}

\paragraph{Provenance and scope.}
This is a fresh proof reconstructed from the original stationary source
formulas and the retained contact substitution. It is not represented as
the original bytes of the former R21 file. Its finite identities are
checked by \texttt{check\_r1\_qo\_compact\_source.py}; the uniform estimates
are proved below. The conclusion concerns the complete scalar on the
two-normal graph. Both state normals and the outer testing charts remain
separate obligations. No historical review receipt is used as a premise.
Source quantities lie in a real closed extension $K_c$ and tests and
states in an independent real closed extension $K_t$; every SOS assertion
is in $K_c\otimes_{\mathbb R}K_t$, without imposing an order on that
tensor ring or identifying its two fields.

The native source formulas are the unchanged formulas in
\texttt{inherited\_r19/inherited\_r18/inherited\_r17/definitions/}
\texttt{weak\_c2\_source\_formulas.tex}. We use their actual common-II
source, with $d=1$, $\lambda=\beta=0$, all original strict gaps and
$n\geq0$. In the original R1 sector $B_s,t_s,a_sB_s,t_sB_s$ are
positive infinitesimals. Write $r,h\in(0,1)$ for its two native source
parameters, $c_s=1+a_sB_sr$, $R_s=1+ha_sB_s$, and $K_s=1+a_sB_s$.
The sector treated here is
\[
 q_{\rm old}=\frac{R_s^2}{a_st_sK_s^2}\quad\hbox{unbounded}.
 \tag{RSR1}
\]
The original positive gap $c_s-(1-h)/t_s$ implies $1-h$ infinitesimal;
therefore $h\sim1$. Since $R_s,c_s,K_s\sim1$, RSR1 is equivalent to
$a_st_s$ infinitesimal. No condition on the size of $a_s$ itself is
introduced. The independent testing variables below are $u,v>0$
infinitesimal and $z>0$ arbitrary.

\subsection{The exact compact source and its uniform bounds}

Define pure-source quantities
\[
 j=\frac{a_sc_s}{R_s},\quad B=\frac{B_sR_s}{c_s},\quad
 T=\frac{t_sc_s}{R_s},\quad p=\frac{a_s(1-r)}{R_s},\quad
 H=\frac{1-h}{t_sc_s},
 \qquad x=jB,\quad y=TB,\quad q_z=jT .
 \tag{RSR2}
\]
The native gaps give $0<p<\min(1,j)$ and $0<H<1$.
Here $B,T,x,y,q_z$ are positive infinitesimals, while $j$ is arbitrary
positive. Direct inversion gives
\[
 r=\frac{j-p}{j+px},\quad h=\frac{1-HT}{1+HTx},\quad
 a_s=\frac{j+px}{1+HTx},\quad
 B_s=\frac{B(1+HTx)}{1+pB},\quad
 t_s=\frac{T(1+pB)}{1+HTx}.
 \tag{RSR3}
\]
In particular the compact variables describe the actual source,
including its $r\to0$, $p\to1$, $H\to0$, and $H\to1$ openings.
Set
\[
 K=j+1-p,\qquad E=1+B(HT+p-1),\qquad
 W_0=(1-H)j+1-p.
 \tag{RSR4}
\]
Then $\tfrac12(j+1)<K<j+1$, $E\sim1$, and $W_0>0$.
The following polynomials contain the entire stationary dictionary:
\begin{align*}
\mathcal P={}&j+1+B\{(j+1)^2-p\}\\
 &+TB\{Hj+j+1+B[Hj(2j+2-p)+p(2j+1-p)]\}\\
 &+HjT^2B^2(1+Bp+BHj),\\
\mathcal L={}&K+T(W_0-H)+BT[pK+Hj(1-p)]+BHT^2(W_0-1),\\
\mathcal N={}&p+B(-j^2+3jp-j-p^2+p)\\
 &+BT(Hj^2+Hj-j^2+jp-j+p)\\
 &+B^2T(2Hj^2p-Hj^2-Hjp^2+2Hjp
             -2j^2p+3jp^2-jp-p^3+p^2)\\
 &+B^2HT^2j(Hj-j+p)(1+Bp).
 \tag{RSR5}
\end{align*}
With $\mu=m(1-p)$, the exact six original stationary weights are
\[
\begin{aligned}
 \mu&=\frac{jp(1+y)(1+x)(1+Hy)}{2\mathcal P},&
 n&=\frac{\mathcal N}{2(j-p)\mathcal P},\\
 C&=\frac{\mathcal L}{2HT\mathcal P},&
 A&=\frac{T(1+x)(1+Hy)(K+Hjy)}{2(1-HT)\mathcal P},\\
 D&=\frac{(1+Hy)(K+Hjy)\{1-y[(1-H)j-p]\}}
             {2T(1-H)\mathcal P},&
 \alpha&=\frac{B^2(K+Hjy)\mathcal L}{2E\mathcal P}.
\end{aligned}
 \tag{RSR6}
\]
For example, the original stationary denominator and its source mass
satisfy
\[
 \Delta=\frac{(1+x)\mathcal P}{(1+Bp)(1+HTx)},\qquad
 M=\frac{(1+x)\mathcal L}{(1+HTx)^2}.
\]
The two native strict gaps become
\[
 g=\frac{(1-p)(1+x)}{1+HTx},\qquad
 j_{\rm gap}=\frac{(1-H)(1+x)}{1+Bp}.
\]
These equations verify the source origin of RSR6 directly; no new
stationary solution has been substituted for the original one.

Every term after $j+1$ in $\mathcal P$, divided by $j+1$, is
infinitesimal: use $B(j+1)=B+x\in\mathrm I$, $T\in\mathrm I$,
$p<1$, and $H<1$. The same termwise estimate, using
$K>(j+1)/2$, applies to $\mathcal L-K$. Thus
\[
 \mathcal P\sim j+1,\qquad \mathcal L\sim K.
 \tag{RSR7}
\]
These are relative equivalences, valid also when $j$ is infinitesimal
or unbounded. They show, in particular, that $A/T$, $THC$, and
$T(1-H)D$ are positive bounded units, $\mu/p$ is bounded, and
$h_0=\mu+1/2$ is a positive bounded unit. Moreover $\alpha/B$ is
infinitesimal. Only the displayed source openings can make $m,n,C,D$
unbounded.

The original physicality condition gives a useful uniform restriction.
Regroup the exact polynomial $\mathcal N$ as
\[
\mathcal N=p\{1+B(3j+1-p)+BT(j+1)\}
 -Bj(j+1)\{1+(1-H)T\}+\mathcal R.
 \tag{RSR8}
\]
Here $\mathcal R$ is precisely the last two lines of $\mathcal N$
in RSR5. Since $0<p<\min(1,j)$, each coefficient in its $B^2T$
line is bounded in absolute value by a fixed multiple of $j(j+1)$.
The $B^2HT^2$ line has the same bound, multiplied by the bounded
factor $T(1+Bp)$. Consequently
\[
 |\mathcal R|\leq C_0BT\,Bj(j+1)
\]
for a fixed real $C_0$. All corrections multiplying $p$ and
$Bj(j+1)$ in RSR8 are infinitesimal. As $n\geq0$ is exactly
$\mathcal N\geq0$, it follows that
\[
 \eta_*:=\frac{Bj(j+1)}p\ \hbox{is bounded},\qquad
 \eta_*\in\mathrm I\ \Longrightarrow\
 n\sim\frac{p}{2(j-p)(j+1)}>0.
 \tag{RSR9}
\]
Indeed $\mathcal N=p(1+\varepsilon_1)-Bj(j+1)(1+\varepsilon_2)$
with $\varepsilon_i\in\mathrm I$. This proves both assertions without
dividing by $n$ and includes literal $n=0$ in the bounded-$\eta_*$
alternative. It also gives $nr$ bounded.

\subsection{The complete graph and its actual contact circuit}

The quarter-parent, before any packet is removed, is
\begin{align*}
\mathscr P={}&m\{vW+uvY-X-uZ\}^2
 +CD_v(uW-Y)^2+AD_uD_vY^2\\
&+(D+\tfrac12)(uW+uvX-Y-vZ)^2
 +(\alpha+\tfrac12)D_u(vY-Z)^2+nD_uD_vX^2\\
&+\tfrac12D_v(X+uZ)^2-\tfrac12D_uD_vZ^2,
\qquad D_u=1+u^2,\ D_v=1+v^2.
 \tag{RSR10}
\end{align*}
The source-free rotated state coordinate is $w=(W+uY)/D_u$.
On its nonzero chart set $w=1$ and $z=(Z-uX)/D_u$. Define
the two complete state normals
\[
 N=X+uZ-pD_uv,\qquad M=uW-Y-HD_uvz.
\]
Their exact inverse is
\[
\begin{aligned}
 X&=N/D_u+pv-uz,& Z&=z+uN/D_u+puv,\\
 Y&=u-M/D_u-Hvz,& W&=1+uM/D_u+uHvz.
\end{aligned}
 \tag{RSR11}
\]
The scalar studied here is exactly
$F_{\rm full}=(\mathscr P/D_u)|_{N=M=0}$.
No $n$ packet, $m$ packet, or state-normal term is silently removed.
Substitution of RSR11 into RSR10 gives the complete graph RB2 of
\texttt{r1\_qo\_contact\_circuit\_budgets.tex}. The independent generic
parent calculation in the checker verifies that equality before
specializing any stationary weight.

Use the separated rows
\[
 U=u/\sqrt x,\quad V=v/\sqrt y,\quad Z_*=z/\sqrt{q_z},\qquad
 Q_*=U-VZ_*,\quad B_*=V-UZ_*,\quad R_*=Z_*-UV,
\]
\[
 \Psi=U^2+V^2+Z_*^2+U^2V^2Z_*^2-4UVZ_*.
 \tag{RSR12}
\]
These are rows; $U,V,Z_*$ are not assumed bounded coefficients.
The actual contact is $(U,V,Z_*;N,M)=(1,1,1;0,0)$.
The exact native formulas make the complete parent value and all five
first derivatives vanish there. On the positive core,
\[
 \Psi=(U-V)^2+Z_*^2(1-UV)^2+2UV(1-Z_*)^2
 \tag{RSR13}
\]
is a legal separated sum of squares: the last root is
$\sqrt2\,(uv)^{1/2}(xy)^{-1/4}(1-Z_*)$, with its radicals in
the two separate real closed fields.

Let the ten positive monomial coefficients and the seven circuit
masses be exactly RB6--RB8. In particular
\[
 M_a=Ax,\quad M_b=\{\mu(1-p)+(n+\tfrac12)p^2\}y,\quad
 M_c=\alpha q_z,\quad M_{ac}=nj^2y,
\]
\[
 f_c=M_b-M_{ac},\quad q_c=M_a-f_c,\quad
 e_c=(n+\tfrac12)p^2y^2-g_c,\quad d_c=M_{b^2c}-k_c.
\]
The original contact balances give the exact scalar identity
\begin{align*}
 F_{\rm full}={}&q_cQ_*^2+M_{ac}B_*^2+r_cR_*^2
  +d_cV^2Q_*^2+e_cV^2B_*^2\\
 &+f_c\Psi+g_cV^2\Psi+k_cZ_*^2(V^2-1)^2.
 \tag{RSR14}
\end{align*}
The identities RB7 and the two odd coefficient balances are independently
checked on the actual source. Thus RSR14 is a finite polynomial identity,
not a Taylor approximation to the contact.

Here is an explicit source-only formula for its radial mass. Put
$\mathcal F=K+T\mathcal F_1+T^2\mathcal F_2$, where
\begin{align*}
\mathcal F_1={}&-2BHjp+BHj+2Bjp-2Bp^2+2Bp\\
 &-Hj-Hp-H+j+1-p^2/j+p/j,\\
\mathcal F_2={}&-B^2H^2jp+B^2Hjp-B^2Hp^2+B^2Hp\\
 &-BH^2j-BH^2p+BHj-BHp^2/j+BHp/j.
 \tag{RSR15}
\end{align*}
Every displayed term divided by $K$ is bounded, using $p/j<1$;
terms containing $Bj$ are even infinitesimal. Hence $\mathcal F\sim K$.
Direct source substitution yields
\[
\begin{aligned}
 f_c&=\frac{Tx^2\mathcal F}{2\mathcal P}>0,&
 q_c&=Ax-f_c\sim\frac{TxK}{2(j+1)}>0,\\
 r_c&=\frac{B^2jT(1-p)(1+x)(1+Hy)\mathcal L}
                {2E\mathcal P}>0,&
 g_c&=\frac{HT^2j^2B^3\mathcal L}{2\mathcal P}>0.
\end{aligned}
 \tag{RSR16}
\]
In particular
\[
 f_c\sim\frac{Tx^2K}{2(j+1)},\qquad
 \frac{r_c}{f_c}\sim\frac{1-p}{j},\qquad
 \frac{g_c}{f_c}=Hy\frac{\mathcal L}{\mathcal F}.
 \tag{RSR17}
\]

For the remaining signed mass define
\begin{align*}
\mathcal K={}&pW_0+
 B(1-p)\{j[H(p+1)-p]-(1-p)(p-H)\}\\
 &+TBH\{(p+1)(1-H)j+(1-p)(BHj+p+1-H)\}\\
 &+T^2B^2H^2j(1-H).
 \tag{RSR18}
\end{align*}
Then
\[
 k_c=-\frac{T^2B^2j(1+x)\mathcal K}{2E\mathcal P}.
\]
The identity
\begin{align*}
\mathcal K-pW_0+Bp(1-p)(j+1)
={}&B(1-p)\{Hj(p+1)+H(1-p)+p^2\}\\
 &+TBH\{(p+1)(1-H)j+(1-p)(BHj+p+1-H)\}\\
 &+T^2B^2H^2j(1-H)
 \tag{RSR19}
\end{align*}
has a nonnegative right side. Since $W_0\geq1-p$ and
$B(j+1)\in\mathrm I$, it proves $\mathcal K>0$ and $k_c<0$.
The displayed finite arrays also give
\[
 \mathcal K/K=O(p+B),\qquad
 |k_c|/(q_cy)=O(p+B).
 \tag{RSR20}
\]
All constants in these estimates are real and uniform over the source
openings; neither $j-p$ nor $1-H$ has been bounded below.

\subsection{Paying the signed masses, with the full current reserved}

Put
\[
 h_0=\mu+\tfrac12,\qquad
 M_e=\frac{nh_0j^2}{n+h_0},\qquad J_c=M_ey,
 \qquad e_0=\tfrac12p^2y^2-g_c.
 \tag{RSR21}
\]
These are pure-source quantities and $0\leq J_c\leq M_{ac}$.
Writing $e_c=np^2y^2+e_0$ preserves the entire second correlated
current packet. The conclusion is the following precise scalar
statement: after subtracting the reserved positive packets
\[
 (M_{ac}-J_c)B_*^2+np^2y^2V^2B_*^2,
 \tag{RSR22}
\]
the remainder is a separated SOS and retains, for example, one half
of $q_cQ_*^2+r_cR_*^2+f_c\Psi+J_cB_*^2$.
When $n=0$ the two reserved packets and $J_c$ vanish literally.
This assertion supplies precisely the effective current mass needed
when the full $N$ normal is later completed against the original $n$
packet. It is not a completion of that normal itself.

Here and below ``infinitesimal cost'' means that the budget multiplies
the original positive packet by a bounded separated Gram coefficient
with zero residue. Finite sums of such costs can be paid while retaining
any fixed fraction smaller than one. Multiplication of the explicit
SOS identities by pure-source nonnegative scalars and pure-testing
squares is legal. The remaining diagonal coefficients are of the form
$1-\varepsilon$, with $\varepsilon$ infinitesimal in the bounded tensor
algebra, so the inherited bounded-Gram completion applies. Its actual
text is the paragraph ``Why a strict bounded Gram suffices'' in
\texttt{inherited\_r19/parent\_reference.tex}; only that field-independent
lemma is imported, not the C14.2 source assumptions in the same file.
For clarity, the elementary step is, for a pure-source infinitesimal
$s$ and bounded pure-testing $t$, to choose real $a,\varepsilon>0$
with $|t|<a$ and $|s|<\varepsilon/a$ and use
\[
 \varepsilon\pm st=\tfrac12\{(\varepsilon/a+s)(a\pm t)
                         +(\varepsilon/a-s)(a\mp t)\}.
\]
Every summand has separated nonnegative factors and hence a separated
square root. Interchange the fields for the other type of infinitesimal
entry and use a finite real diagonal budget. This proves the exact
SOS lifting step used here. No square root or inverse of a mixed sum
is used.

First $e_0^-\leq g_c$. The exact budget
\[
 (1+V^2)\Psi-V^2B_*^2
 =Q_*^2+R_*^2+V^2(1-UVZ_*)^2
 \tag{RSR23}
\]
pays its negative part from the circuit at cost
\[
 \frac{e_0^-}{f_c}+\frac{e_0^-}{f_cy}v^2
 \ \leq\ Hy\frac{\mathcal L}{\mathcal F}
       +H\frac{\mathcal L}{\mathcal F}v^2\in\mathrm I.
\]
Its positive part is retained as a square. Also $d_c>0$ because
$M_{b^2c}>0$ and $k_c<0$; the $d_c$ and $g_c$ terms cause no loss.

For $k_c$, use the exact identity
\[
 Z_*(V^2-1)=-VQ_*-R_*.
 \tag{RSR24}
\]
The $VQ_*$ part is paid from $q_cQ_*^2$ with coefficient
$2|k_c|v^2/(q_cy)$, which is infinitesimal by RSR20. The remaining
$R_*$ part is paid in the following exhaustive three cases. Fixed
real multipliers such as the factor two do not change any cost claim.

\paragraph{Case 1: $1-p$ has a positive real lower bound.}
The exact ratio is
\[
 \frac{|k_c|}{r_c}
 =\frac{T\mathcal K}{(1-p)(1+Hy)\mathcal L}\in\mathrm I.
 \tag{RSR25}
\]
This uses only $\mathcal K/\mathcal L=O(p+B)$ and $T\in\mathrm I$.
Neither current mass nor a bound on $j$ is required.

\paragraph{Case 2: $1-p\in\mathrm I$ and $\eta_*$ has a positive real
lower bound.}
Here $p\sim1$, so $Bj(j+1)$ is a bounded positive unit. Since $B$ is
infinitesimal, $j$ is unbounded and $j^2B$ is a positive bounded unit.
Equations RSR16--RSR20 give
\[
 \frac{|k_c|}{f_c}\in\mathrm I,\qquad
 \frac{|k_c|}{f_cx}
 =O\!\left(\frac{T(p+B)}{j^2B}\right)\in\mathrm I.
 \tag{RSR26}
\]
The exact circuit budget
\[
 (1+U^2)\Psi=R_*^2+B_*^2+U^2Q_*^2+U^2(1-UVZ_*)^2
 \tag{RSR27}
\]
therefore pays $|k_c|R_*^2$ at cost
$|k_c|/f_c+(|k_c|/(f_cx))u^2\in\mathrm I$.
This case includes any intersection with literal $n=0$.

\paragraph{Case 3: $1-p,\eta_*\in\mathrm I$.}
Now $p\sim1$, $j>p$, and RSR9 gives $n>0$ together with its exact
relative estimate. The source mass $M_e$ is a positive bounded unit.
For the lower bound use
\[
 \frac{p^2}{M_e}
 =\frac{p^2}{h_0j^2}+\frac{p^2}{nj^2}=O(1).
\]
For the upper bound use $M_e\leq h_0j^2$ when $j$ is bounded,
and $M_e\leq nj^2\sim p/2$ when $j$ is unbounded. In particular,
\[
 \frac{|k_c|}{f_c}\in\mathrm I,\qquad
 \frac{|k_c|}{q_cy}=O(1),\qquad
 \frac{|k_c|}{J_cx}
 =O\!\left(\frac{T(p+B)(1+1/n)}{j^2}\right)=O(1).
 \tag{RSR28}
\]
The last bound follows directly from
$1/n\sim2(j-p)(j+1)/p$, with $p\sim1$ and $j>p$.
There is an exact separated SOS budget
\[
 2\{\Psi+V^2Q_*^2+U^2B_*^2\}-R_*^2\quad\hbox{is SOS}.
 \tag{RSR29}
\]
For completeness, set $t_*=UV$, $\delta_*=U-V$,
$E_*=Z_*(1-t_*)$, $T_*=t_*(1-Z_*)$, and
$J_*=T_*-\delta_*^2Z_*$. Its SOS decomposition is
\[
 (E_*+T_*)^2+2\delta_*^2+4t_*(1-Z_*)^2
                       +2J_*^2+4t_*\delta_*^2Z_*^2.
\]
All roots are separated as in RSR13. Applying RSR29 costs a multiple
of $|k_c|/f_c$ on the circuit, $|k_c|v^2/(q_cy)$ on $Q_*$,
and $|k_c|u^2/(J_cx)$ on $B_*$. Each is infinitesimal by RSR28.

The three cases exhaust the source because $0<1-p<1$ and $\eta_*$
is bounded. A finite number of the displayed allocations therefore
proves the scalar assertion following RSR22 with a strict retained
fraction. All coefficient bounds depend only on the source and
$u,v$; no bound on the independent state coordinate $z$ was used.
The proof allows arbitrary relative rates of all source and testing
infinitesimals. Literal $n=0$ is covered by Cases 1 and 2; Case 3
excludes it by the original physicality equation itself.

\paragraph{What remains outside this lemma.}
The scalar identity is exact for the complete graph of RSR10. A
full-parent local theorem must also allocate both first-normal arrays
and the complete quadratic block from RSR11, using the reserved
current before applying signed scalar costs. The outer testing charts,
the state chart $w=0$, the sign exits and complex-state splitting must
be supplied by their own exact arguments. This file makes none of
those additional admissions by implication.
